% ============================================================
%  Dédicace — Mémoire de Master BIBDA — CHOUFA Zouhair
% ============================================================
\clearpage

% \thispagestyle{empty}
\begin{tikzpicture}[remember picture, overlay]
    % Cadre décoratif discret aux quatre coins de la page
    \draw[line width=1pt, titleblue]
        ([shift={(1.5cm,-1.5cm)}]current page.north west) -- ++(1.4cm,0);
    \draw[line width=1pt, titleblue]
        ([shift={(1.5cm,-1.5cm)}]current page.north west) -- ++(0,-1.4cm);

    \draw[line width=1pt, titleblue]
        ([shift={(-1.5cm,-1.5cm)}]current page.north east) -- ++(-1.4cm,0);
    \draw[line width=1pt, titleblue]
        ([shift={(-1.5cm,-1.5cm)}]current page.north east) -- ++(0,-1.4cm);

    \draw[line width=1pt, bordergreen]
        ([shift={(1.5cm,1.5cm)}]current page.south west) -- ++(1.4cm,0);
    \draw[line width=1pt, bordergreen]
        ([shift={(1.5cm,1.5cm)}]current page.south west) -- ++(0,1.4cm);

    \draw[line width=1pt, bordergreen]
        ([shift={(-1.5cm,1.5cm)}]current page.south east) -- ++(-1.4cm,0);
    \draw[line width=1pt, bordergreen]
        ([shift={(-1.5cm,1.5cm)}]current page.south east) -- ++(0,1.4cm);

    % Petit motif de noeud sémantique (rappel Knowledge Graph) en filigrane
    \begin{scope}[opacity=0.55, shift={($(current page.center)+(0,-8.5cm)$)}]
        \node[circle, draw=lineblue, fill=lineblue!15, minimum size=5pt, inner sep=0pt] (n1) at (0,0) {};
        \node[circle, draw=titleblue, fill=titleblue!15, minimum size=5pt, inner sep=0pt] (n2) at (0.9,0.35) {};
        \node[circle, draw=bordergreen, fill=bordergreen!15, minimum size=5pt, inner sep=0pt] (n3) at (1.6,-0.25) {};
        \node[circle, draw=lineblue, fill=lineblue!15, minimum size=5pt, inner sep=0pt] (n4) at (-0.9,0.3) {};
        \node[circle, draw=titleblue, fill=titleblue!15, minimum size=5pt, inner sep=0pt] (n5) at (-1.6,-0.2) {};
        \draw[lineblue!70] (n1) -- (n2);
        \draw[titleblue!70] (n2) -- (n3);
        \draw[bordergreen!70] (n1) -- (n3);
        \draw[lineblue!70] (n1) -- (n4);
        \draw[titleblue!70] (n4) -- (n5);
        \draw[bordergreen!70] (n1) -- (n5);
    \end{scope}
\end{tikzpicture}

\vspace*{2.2cm}
\begin{center}
    {\Large \scshape \color{titleblue} \textbf{Dédicace}} \\[0.15cm]
    \textcolor{bordergreen}{\rule{4.5cm}{0.8pt}}
\end{center}

\vspace{1.8cm}

\begin{center}
\begin{minipage}{0.82\textwidth}
\setlength{\parindent}{0pt}
\raggedright
\Large
\itshape

À mes très chers parents, \\[0.3cm]
\normalsize\upshape
dont l'amour inconditionnel, les sacrifices silencieux et la confiance
jamais démentie ont tracé, bien avant moi, le chemin de chacune de mes
réussites. Que ce modeste travail soit le reflet, même imparfait, de
tout ce que je vous dois.

\vspace{0.9cm}
\Large\itshape
À mes frères et sœurs, \\[0.3cm]
\normalsize\upshape
pour votre soutien constant, vos encouragements dans les moments de
doute, et la complicité qui nous unit depuis toujours.

\vspace{0.9cm}
\Large\itshape
À mes amis et collègues de la promotion BIBDA, \\[0.3cm]
\normalsize\upshape
compagnons de cette aventure académique exigeante, avec qui chaque
nuit blanche de débogage et chaque réussite ont pris un sens partagé.

\vspace{0.9cm}
\Large\itshape
À tous ceux qui croient qu'une donnée brute, aussi hétérogène
soit-elle, peut devenir connaissance structurée pour peu qu'on lui
consacre la rigueur et la patience nécessaires, \\[0.3cm]
\normalsize\upshape
ce mémoire vous est dédié.

\end{minipage}
\end{center}

\vfill
\begin{center}
    \textcolor{lineblue}{\rule{6cm}{0.6pt}}\\[0.25cm]
    {\small \itshape ``Les données ne parlent pas d'elles-mêmes ;\\
    c'est la structure qu'on leur donne qui les fait parler.''}
\end{center}
\vspace{1.5cm}

\newpage